\section{Bounded residues and rationality}\label{sec:family}
The residue theorem uses the same separation principle: force neighbouring
coefficients to vanish while retaining one central value.

\begin{theorem}\label{res:residueseries}
For every $m\ge3$,
\[
 A_m=\sum_{n\ge1}\frac{\ph(n)\bmod m}{2^n}\notin\Q.
\]
For $k\ge1$ and $f:\Z/2^k\Z\to\Q$, the series
$\sum_{n\ge1}f(\ph(n)\bmod2^k)2^{-n}$ is rational exactly when $f$ is
constant on the even residue classes. If that constant is $c$, its value is
$3f(1)/4+c/4$.
\end{theorem}

\begin{lemma}[A bounded isolated pulse]\label{lem:bounded-pulse}
Let $a_n\in\Z$ satisfy $|a_n|\le C$. Suppose that for arbitrarily large $L$
there is $N>L$ such that $a_N\ne0$ and
$a_{N+t}=0$ for $0<|t|\le L$. Then $\sum_{n\ge1}a_n2^{-n}$ is irrational.
\end{lemma}
\begin{proof}
If the sum has denominator $q$, every scaled tail
$T_j=\sum_{t\ge1}a_{j+t}2^{-t}$ belongs to $q^{-1}\Z$.
Choose $L$ so large that $C2^{-L}<1/q$.
The right zero block gives $|T_N|\le C2^{-L}<1/q$, hence $T_N=0$.
The left zero block now gives
$T_{N-L-1}=a_N2^{-L-1}$, a nonzero element of $q^{-1}\Z$ with absolute
value less than $1/q$, a contradiction.
\end{proof}

\begin{proof}[Proof of Theorem~\ref{res:residueseries}]
Fix a unit $s\pmod m$ and a block length $L$. For each nonzero
$t\in[-L,L]$, choose a distinct prime $q_t\equiv1\pmod m$ larger than $L$
and outside the prime divisors of $m$. Prescribe
\[
 p\equiv s\pmod m,\qquad p\equiv-t\pmod{q_t}.
\]
This is a reduced progression, because each $q_t>|t|$. Dirichlet supplies
arbitrarily large primes $p$ in it. Then $m\mid\ph(p+t)$ for every
$0<|t|\le L$, whereas $\ph(p)\equiv s-1\pmod m$.

Subtract $f(0)$ from the observable and clear its denominators. Whenever
$f(s-1)\ne f(0)$, the preceding construction gives arbitrarily long bounded
isolated pulses, so Lemma~\ref{lem:bounded-pulse} applies. For the
least-residue map choose $s=-1$: its central residue is $m-2\ne0$ when
$m\ge3$. For $m=2^k$, every even residue $r$ has $r+1$ a unit, so a
nonconstant restriction to the even residues supplies such a pulse.
Conversely, $\ph(n)$ is even for every $n\ge3$. Constancy on the even
classes therefore makes all terms after $n=2$ constant, giving the stated
rational value. This converse requires no recurrent-support theorem.
\end{proof}
Fixed bounded observables do not recover the unbounded coefficients of
$S$. In the pulse proof, the coefficient bound is fixed before the zero
block is chosen; allowing the modulus to grow changes that bound.
